% ============================== CI-5 ==============================
\reading{CI-5}{Monetary and Fiscal Policy: Safeguarding Stability and Trust}{STRIKE FIRST}{2}

\begin{cikeybox}[title={The three objectives, in the spine's own words}]
\begin{enumerate}[label=\textbf{\alph*.}]
\item Compare and contrast the channels through which fiscal policy and monetary policy
influence a country's economic activity and financial markets, and define the ``region of
stability'' in terms of their joint policy stances.
\item Describe the consequences of breaching the boundaries of the region of stability, and
how these consequences have evolved over time in advanced economies and in emerging market
economies.
\item Describe the risks that global economies face as a result of high public debt levels,
including the potential for these high debt levels, in combination with other factors, to
drive tension between fiscal policy and monetary policy.
\end{enumerate}
\greyline{Source: Annual Economic Report 2023, Section 2, BIS (June 2023). The class notes
label this Reading 104; it is CI-5 on the 2026 spine. The notes' heading for objective (c)
reads ``iscal policy'', a dropped ligature; the heading above is the spine's wording.}
\end{cikeybox}

% ------------------------------------------------------------------
\lo{CI-5 a}{Compare and contrast the channels through which fiscal policy and monetary policy influence a country's economic activity and financial markets, and define the ``region of stability'' in terms of their joint policy stances.}

\subsection*{Channels of policy influence}

\begin{center}
\renewcommand{\arraystretch}{1.25}
\begin{tabularx}{\textwidth}{@{}>{\raggedright\arraybackslash}X >{\raggedright\arraybackslash}X@{}}
\toprule
\rowcolor{PaleNavy}
{\sffamily\bfseries\color{FRMNavy}Monetary policy} &
{\sffamily\bfseries\color{FRMGold}Fiscal policy}\\
\midrule
Managed by central banks & Managed by governments\\
Influences the economy \term{indirectly}, via interest rates and money supply &
Influences the economy \term{directly}, via taxes, spending, and transfer payments\\
Tools: open market operations, policy interest rates &
Tools: taxation, public spending, subsidies\\
Impacts: bond yields, borrowing costs, asset prices, exchange rates &
Impacts: aggregate demand, employment, income distribution\\
\bottomrule
\end{tabularx}
\end{center}
\figcap{Table CI-5.1. Row two is the whole contrast in one word each. Monetary policy reaches
the economy \emph{indirectly} through the price of money, fiscal policy reaches it
\emph{directly} through the government's own spending and taxing. Everything else in the
table follows from that difference, including why the impacts row lists market prices on one
side and real-economy quantities on the other.}

\begin{cifigblock}
{\centering
\begin{tikzpicture}[
  font=\sffamily\scriptsize,
  gl/.style={draw=FRMGold, fill=PaleGold, text width=3.6cm, align=center,
             minimum height=0.62cm, inner sep=3pt, rounded corners=1pt,
             font=\sffamily\bfseries\scriptsize},
  st/.style={text width=4.3cm, align=center, inner sep=2pt},
  hd/.style={draw=FRMRule, fill=FRMSoft, text width=4.6cm, align=center,
             minimum height=0.6cm, inner sep=3pt, rounded corners=1pt,
             font=\sffamily\bfseries\scriptsize\color{FRMNavy}},
  dn/.style={-{Stealth[length=4pt]}, FRMNavy, line width=0.7pt}
]
\node[gl] (goal) at (4.6,0) {Goal: reduce inflation};

\node[hd] (mh) at (0,-1.05) {Monetary policy};
\node[hd] (fh) at (9.2,-1.05) {Fiscal policy};
\draw[dn] (goal.south) -- ++(0,-0.18) -| (mh.north);
\draw[dn] (goal.south) -- ++(0,-0.18) -| (fh.north);

\node[st, below=4mm of mh] (m1) {\textbf{a)} policy rate $\uparrow$};
\node[st, below=3.5mm of m1] (m2) {loans and borrowing costlier $\downarrow$};
\node[st, below=3.5mm of m2] (m3) {consumer spending $\downarrow$};
\node[st, below=3.5mm of m3] (m4) {demand $\downarrow$, prices $\downarrow$};
\node[st, below=3.5mm of m4, font=\sffamily\bfseries\scriptsize\color{FRMGreen}] (m5)
  {inflation $\downarrow$};
\foreach \a/\b in {m1/m2, m2/m3, m3/m4, m4/m5} \draw[dn] (\a) -- (\b);

\node[st, below=9mm of m5] (n1) {\textbf{b)} open market operations};
\node[st, below=3.5mm of n1] (n2) {central bank \emph{sells} bonds to the public};
\node[st, below=3.5mm of n2] (n3) {money is taken from investors, so money supply
  $\downarrow$};
\node[st, below=3.5mm of n3, font=\sffamily\bfseries\scriptsize\color{FRMGreen}] (n4)
  {inflation $\downarrow$};
\foreach \a/\b in {n1/n2, n2/n3, n3/n4} \draw[dn] (\a) -- (\b);

\node[st, below=4mm of fh] (f1) {\textbf{a)} government spending $\downarrow$};
\node[st, below=3.5mm of f1] (f2) {infrastructure $\downarrow$, subsidies $\downarrow$};
\node[st, below=3.5mm of f2] (f3) {money supply $\downarrow$};
\node[st, below=3.5mm of f3, font=\sffamily\bfseries\scriptsize\color{FRMGreen}] (f4)
  {inflation $\downarrow$};
\foreach \a/\b in {f1/f2, f2/f3, f3/f4} \draw[dn] (\a) -- (\b);

\node[st, below=9mm of f4] (g1) {\textbf{b)} taxes $\uparrow$};
\node[st, below=3.5mm of g1] (g2) {less money with people};
\node[st, below=3.5mm of g2] (g3) {demand $\downarrow$, prices $\downarrow$};
\node[st, below=3.5mm of g3, font=\sffamily\bfseries\scriptsize\color{FRMGreen}] (g4)
  {inflation $\downarrow$};
\foreach \a/\b in {g1/g2, g2/g3, g3/g4} \draw[dn] (\a) -- (\b);

\draw[FRMRule, line width=0.8pt] (4.6,-1.45) -- (4.6,0 |- n4.south);
\end{tikzpicture}\par}
\figcap{Figure CI-5.1. The same goal reached down four different chains, redrawn from the
notes' hand-written working. Read the second step of each chain: the two monetary routes act
on the \emph{price} or the \emph{quantity} of money, while the two fiscal routes act on the
government's own spending or on households' after-tax income. Note the direction in the open
market operations chain, which is the one that reverses most often under pressure: to reduce
inflation the central bank \emph{sells} bonds, which drains money from the public.}
\end{cifigblock}

\subsection*{Interaction of monetary and fiscal policy}

\begin{enumerate}
\item High fiscal spending $\rightarrow$ increased demand $\rightarrow$ inflation
$\rightarrow$ central bank responds with \term{rate hikes} (tight monetary policy).
\item Tight monetary policy = higher interest rates $\rightarrow$ \term{costlier government
borrowing} (on new debt and refinancing).
\end{enumerate}

This creates a \term{policy feedback loop}, where fiscal expansion triggers monetary
tightening, and monetary tightening strains fiscal budgets.

\subsection*{The region of stability}

\begin{cidefbox}[title={Region of stability}]
A balanced zone where both policies support stable, healthy growth. Hence, a \term{margin of
safety} is essential to avoid pushing fiscal or monetary policy too far.
\end{cidefbox}

\begin{enumerate}[label=\textbf{\alph*)}]
\item The region of stability is country-specific, hard to define, and constantly shifting.
\item It is influenced by gradual changes (like demographics) and sudden shocks (like
geopolitical events).
\item Policies like regulation or FX intervention can alter its boundaries.
\end{enumerate}

\begin{cifigblock}
{\centering
\begin{tikzpicture}[font=\sffamily\scriptsize]
% axes
\draw[-{Stealth[length=5pt]}, FRMNavy, line width=0.8pt] (0,0) -- (10.9,0)
  node[right, font=\sffamily\scriptsize\color{FRMNavy}] {};
\draw[-{Stealth[length=5pt]}, FRMNavy, line width=0.8pt] (0,0) -- (0,6.2);
\node[below=16pt, font=\sffamily\bfseries\scriptsize\color{FRMNavy}] at (4.0,0)
  {FISCAL STANCE \ \ \textnormal{tight $\longrightarrow$ loose}};
\node[rotate=90, above=10pt, font=\sffamily\bfseries\scriptsize\color{FRMNavy}] at (0,3.1)
  {MONETARY STANCE \ \ \textnormal{tight $\longrightarrow$ loose}};

% AE region
\fill[PaleGreen, draw=FRMGreen, line width=0.9pt, rounded corners=14pt, opacity=0.95]
  (1.0,1.0) rectangle (6.0,4.5);
\node[font=\sffamily\bfseries\scriptsize\color{FRMGreen}, align=center,
      fill=PaleGreen, inner sep=2pt] at (3.5,1.42)
  {region of stability,\\advanced economy};

% EME region
\fill[PaleGold, draw=FRMGold, line width=0.9pt, rounded corners=10pt, opacity=0.95]
  (1.6,2.15) rectangle (4.1,3.9);
\node[font=\sffamily\bfseries\scriptsize\color{FRMGold}, align=center] at (2.85,3.0)
  {smaller region\\for an EME};

% margin of safety
\draw[<->, FRMRed, line width=0.8pt] (6.0,2.6) -- (7.4,2.6);
\node[font=\sffamily\scriptsize\color{FRMRed}, fill=white, inner sep=1.5pt, align=center]
  at (6.72,2.15) {margin of safety};

% breach consequences
\node[draw=FRMRule, fill=PaleRed, text width=3.1cm, align=left, inner sep=4pt,
      rounded corners=1pt] (br) at (9.05,4.6)
  {\textbf{Breaching the boundary}
   \begin{itemize}[leftmargin=9pt, topsep=2pt, itemsep=1pt]
   \item currency depreciation and capital outflows amplify
   \item bond repricing can trigger systemic risk
   \end{itemize}};
\draw[-{Stealth[length=4pt]}, FRMRed, line width=0.8pt] (br.west) -- (6.15,4.2);

% shifting boundary
\draw[FRMGrey, dashed, line width=0.7pt, rounded corners=14pt]
  (1.7,1.95) rectangle (6.7,5.3);
\node[font=\sffamily\scriptsize\itshape\color{FRMGrey}, fill=white, inner sep=1.5pt,
      align=center] at (3.0,0.45)
  {dashed: the same region at another moment, since the boundary constantly shifts};
\end{tikzpicture}\par}
\figcap{Figure CI-5.2. The region of stability is defined over the \emph{joint} stance of the
two policies, which is why it is drawn as an area rather than a pair of thresholds: a fiscal
stance that is safe alongside tight money may sit outside the region alongside loose money.
Three things the notes state are drawn here. The margin of safety is the distance you keep
from the edge, not the edge itself. The EME region is smaller, which is why the same policy
move can be inside the region for an advanced economy and outside it for an emerging one.
And the dashed outline is the same region at a different moment, since demographics, shocks,
regulation and FX intervention all move the boundary. The shape and position are
illustrative, not measured.}
\end{cifigblock}

% ------------------------------------------------------------------
\lo{CI-5 b}{Describe the consequences of breaching the boundaries of the region of stability, and how these consequences have evolved over time in advanced economies and in emerging market economies.}

Among advanced economies (AEs), the pre-pandemic period can be subdivided into two phases.

\begin{cifigblock}
{\centering
\begin{tikzpicture}[
  font=\sffamily\scriptsize,
  ph/.style={draw=FRMRule, fill=PaleNavy, text width=3.35cm, align=center,
             minimum height=0.62cm, inner sep=3pt, rounded corners=1pt,
             font=\sffamily\bfseries\scriptsize\color{FRMNavy}},
  st/.style={text width=3.5cm, align=center, inner sep=2pt},
  end/.style={text width=3.5cm, align=center, inner sep=3pt,
              font=\sffamily\bfseries\scriptsize\color{FRMRed}},
  dn/.style={-{Stealth[length=4pt]}, FRMNavy, line width=0.7pt}
]
\node[ph] (p1) at (0,0) {Phase 1 \ (1980 to mid-1980s)};
\node[ph] (p2) at (4.2,0) {Phase 2 \ (mid-1980s to 2020)};
\node[ph, fill=PaleGold] (p3) at (8.4,0) {Post-COVID};

\node[st, below=3.5mm of p1] (a1) {tried to fine-tune growth with both fiscal and monetary
  policy};
\node[st, below=3.5mm of a1] (a2) {policy overreach: high inflation and debt};
\node[st, below=3.5mm of a2] (a3) {aggressive rate hikes, shift to price stability};
\node[end, below=3.5mm of a3] (a4) {price instability};
\foreach \x/\y in {p1/a1, a1/a2, a2/a3, a3/a4} \draw[dn] (\x) -- (\y);

\node[st, below=3.5mm of p2] (b1) {globalisation and financial liberalisation};
\node[st, below=3.5mm of b1] (b2) {inflation dropped, rates low};
\node[st, below=3.5mm of b2] (b3) {focus moved to financial imbalances: credit excess,
  rising debt};
\node[end, below=3.5mm of b3] (b4) {financial instability};
\foreach \x/\y in {p2/b1, b1/b2, b2/b3, b3/b4} \draw[dn] (\x) -- (\y);

\node[st, below=3.5mm of p3] (c1) {massive stimulus spent};
\node[st, below=3.5mm of c1] (c2) {rapid recovery};
\node[st, below=3.5mm of c2] (c3) {high inflation, high debt, reduced policy space};
\node[end, below=3.5mm of c3] (c4) {financial vulnerability};
\foreach \x/\y in {p3/c1, c1/c2, c2/c3, c3/c4} \draw[dn] (\x) -- (\y);
\end{tikzpicture}\par}
\figcap{Figure CI-5.3. Three episodes, three different ways of leaving the region. The
end-state row is the one to memorise, because the three are not the same failure: Phase 1
ends in \emph{price} instability, Phase 2 in \emph{financial} instability, and the post-COVID
period in \emph{financial vulnerability} with the policy space already spent. Phase 2 is the
counter-intuitive one, since low inflation and low rates were the conditions under which the
imbalances built.}
\end{cifigblock}

A few factors contribute to the trajectory of monetary and fiscal stimulus.

\begin{itemize}
\item First, when sovereign debt levels are high, fiscal stimulus has \emph{less} stimulative
effect, and more extreme monetary efforts are needed to have an impact.
\item Second, there is a self-reinforcing interaction between fiscal and monetary policy. Due
to foreign exchange rate dependency, emerging market economies (EMEs) are more sensitive to
global financial conditions than AEs.
\end{itemize}

\begin{cikeybox}[title={Implications of this foreign exchange rate dependency}]
\begin{enumerate}
\item The region of stability is \term{smaller} for EMEs than AEs.
\item Boundary breaches amplify currency depreciation and capital outflows.
\item Monetary policy in AEs can have a ripple effect in EMEs.
\item EMEs are shifting to flexible exchange rate regimes.
\end{enumerate}
\end{cikeybox}

% ------------------------------------------------------------------
\lo{CI-5 c}{Describe the risks that global economies face as a result of high public debt levels, including the potential for these high debt levels, in combination with other factors, to drive tension between fiscal policy and monetary policy.}

\subsection*{Risks from high public debt and policy tensions}

High public debt poses serious challenges to both monetary and fiscal policy. It limits
flexibility, increases financial vulnerability, and can cause tensions between inflation
control and economic support goals. This debt strains government finances, complicates the
fight against inflation, and introduces substantial financial stability risks.

\subsection*{1. Intersection of multiple fiscal and economic factors}

Government debt, spending, inflation, and interest rates are closely linked. Inflation has
temporarily boosted tax revenues, making fiscal positions look healthier than they are.

\term{Consequences of high inflation}: if inflation stays high, central banks must keep
interest rates elevated, which raises borrowing costs.

\begin{enumerate}[label=\textbf{\alph*)}]
\item Central bank bond buying post-GFC shortened debt maturity, cutting costs in the short
term but lowering remittances.
\item High fiscal debt increases annual debt servicing costs (expected to reach \term{6\% of
GDP}), and if economic growth slows, due to factors like ageing populations or weaker growth
in emerging markets, it could worsen the debt problem.
\end{enumerate}

\term{Why cutting fiscal spending is difficult.} Inflation, ageing populations, and global
challenges push government spending higher.

\subsection*{2. Countercyclical effects}

Monetary and fiscal policies can work against each other. Large fiscal deficits stimulate
demand but raise inflation, which leads to higher interest rates. These higher rates, in turn,
can weaken economic growth, undermining fiscal efforts. This tension is especially visible in
Emerging Market Economies (EMEs), where weak exchange rates worsen inflation.

\begin{cifigblock}
{\centering
\begin{tikzpicture}[
  font=\sffamily\scriptsize,
  nd/.style={draw=FRMRule, fill=PaleRed, text width=2.6cm, align=center,
             minimum height=0.95cm, inner sep=3pt, rounded corners=1pt},
  fd/.style={draw=FRMRule, fill=PaleNavy, text width=3.9cm, align=left,
             inner sep=4pt, rounded corners=1pt},
  lp/.style={-{Stealth[length=4pt]}, FRMRed, line width=0.8pt}
]
\node[nd] (d1) at (0,1.35)   {large fiscal deficit: spending over taxes};
\node[nd] (d2) at (3.3,1.35) {debt $\uparrow$, money supply $\uparrow$, demand $\uparrow$};
\node[nd] (d3) at (3.3,-0.55){inflation $\uparrow$, so the central bank tightens};
\node[nd] (d4) at (0,-0.55)  {higher rates weaken growth, undermining fiscal efforts};
\draw[lp] (d1) -- (d2);
\draw[lp] (d2) -- (d3);
\draw[lp] (d3) -- (d4);
\draw[lp] (d4) -- (d1);

\node[fd] (fdom) at (8.4,0.9)
  {\textbf{Fiscal dominance} exists when monetary policy is constrained by fiscally imposed
   constraints. It arises from:};
\node[fd, fill=PaleGold, below=2.5mm of fdom] (r1)
  {\textbf{1. Political spending pressure.} Pressure to maintain or increase spending in
   support of fiscal priorities};
\node[fd, fill=PaleGold, below=2.5mm of r1] (r2)
  {\textbf{2. Sovereign debt implications.} A rate increase could spark a sovereign debt
   crisis under certain circumstances};
\end{tikzpicture}\par}
\figcap{Figure CI-5.4. The loop on the left is why the two policies can undo each other:
each arm's remedy is the next arm's problem, and the circuit closes. Fiscal dominance, on the
right, is the state in which the loop stops turning because the central bank can no longer
take the action the loop calls for. The tension is sharpest in EMEs, where a weak exchange
rate worsens the inflation step.}
\end{cifigblock}

\subsection*{3. Financial stability risks}

Sovereign debt is vital for financial stability. It acts as a liquid store of value,
collateral, and a benchmark for pricing other assets. Sovereign credit stress (like downgrades
or defaults) can disrupt risk pricing across the economy.

Governments act as backstops, with central banks serving as lenders of last resort. This
creates a \term{feedback loop}: banks rely on sovereign bonds, and sovereigns rely on bank
stability, especially in EMEs.

High debt raises the risk of sharp bond yield movements, especially during monetary
tightening. Bonds issued at low rates pre-pandemic are now more volatile. In panic phases,
investors may flee sovereign debt for cash, causing liquidity stress.

Loose fiscal and monetary policy increases fragility. If markets sense a breach in the region
of stability, bond repricing can trigger systemic risk.

\subsection*{4. Policy implications}

\begin{enumerate}[label=\textbf{\alph*)}]
\item Monetary policy should focus on price stability, while fiscal policy must remain
conservative to preserve crisis flexibility and avoid fuelling inflation.
\item Central banks should reduce balance sheets gradually to retain future policy space and
avoid distorting markets.
\item Fiscal spending should prioritise long-term growth (like infrastructure), not short-term
transfers.
\item Central bank independence is crucial to avoid political interference.
\end{enumerate}
